\chapter{Status bars and gutters}
\label{ch:statusbar}

\begin{aside}
A status bar is one row of fixed-position information at the
edge of the screen. Mode, cursor position, time, current
file. The hardest part is deciding what not to show.
\end{aside}

\section{Layout}

A status bar is a row docked to the bottom (or top):

\begin{lstlisting}
auto root = column_flex(
    main_content,
    status_bar() | height(1)
);
\end{lstlisting}

The main content stretches; the bar takes one row.

\section{Sections}

A status bar typically has left, centre, and right sections.
Left and right grow to push centre to the middle.

\section{Content choices}

Common entries: mode (insert/normal/visual), file path,
cursor position, modified indicator, encoding, line ending,
time, battery, network.

\section{Truncation}

When the bar is narrower than its content, sections truncate
from the centre or the longest section first. \maya{}'s bar
uses priority: highest-priority entries stay; lowest get
ellipsed first.

\section{Gutters}

A gutter is a vertical strip beside the editor: line
numbers, git change markers, diagnostic icons, breakpoint
dots.

\maya{}'s editor widget has a configurable gutter.

\section{Mode colour}

A common pattern: change the status bar background based on
mode. Vim does this. The cell-diff handles updates
efficiently.

\section{Recap}

\begin{recap}
\begin{itemize}[leftmargin=*]
  \item Bar docks to top or bottom, 1 row tall.
  \item Left/centre/right sections.
  \item Priority-based truncation.
  \item Gutter is the vertical analogue.
\end{itemize}
\end{recap}

\section*{Workshop}

\begin{enumerate}[leftmargin=*]
  \item Build a status bar with mode, path, cursor.
  \item Add priority-based truncation.
  \item Add a left gutter with line numbers and git
    markers.
\end{enumerate}
